\makeatletter \def\input@path{{../../}} \makeatother
\documentclass{article}
\usepackage[utf8]{inputenc}
\usepackage{amsmath, amsthm, amssymb, mathpazo, isomath, mathtools}
\usepackage{subcaption,graphicx,pgfplots}
\usepackage{fullpage}
\usepackage{booktabs}
\usepackage{hyperref}
\usepackage{algorithm, algorithmic}
\usepackage{mathtools}
\usepackage{todonotes}
\usepackage{tcolorbox}
\usepackage{totcount,xfp}
\newtotcounter{totalpts}
%\setcounter{totalpts}{0} % Contador de puntos en homework

\newcommand{\getTotal}{\number\totvalue{totalpts} }
\newcommand{\instruccion}{\begin{tcolorbox}[colback=blue!5!white,colframe=blue!75!black!70!white,title=Instructions]
The homework must be developed by one or two people in a .pdf report submitted through the designated drop box on the Canvas platform, where you must also show the developed code if applicable. For your convenience, you may submit the homeworks in a Jupyter Notebook, so it is more comfortable to show the code. In any case, a single file must be submitted as the answer to the homework. The late policy will be: a linear factor will be calculated that is 1 at the time of delivery and 0 48 hours later. This will multiply your obtained score.


The homework has \getTotal points. In each question, the given score will be: the total if it is correct, half if it has a conclusive development but contains errors, and 0.0 points if it has little to no development.

The use of language models (like ChatGPT) is allowed to support the development of the homework, although we suggest trying each question a bit first to generate doubts. While we cannot control if you ask the model directly for the answer, we believe it is better for your learning process to ask for suggestions rather than solutions. For this reason, if you decide to use it, you must attach to your homework another pdf showing the prompts used. We suggest using prompts of this style: "I am a student of [Major] and I am developing the following exercise for a homework in the course [Course]: [Copy exercise]
  And we have seen so far the contents in the attached slides [attach pdfs of the lectures]. Given this content, I would like to ask you for a suggestion to solve [explain difficulty], since I am confused about [write doubt]. As an answer, I do not want you to give me the solution, but a guide to help me start developing my answer."

  If you have doubts and/or comments, please contact us by email or Canvas. It is possible that the formulation has errors.
\end{tcolorbox}}


\newcommand{\code}[1]{\texttt{#1}}
% Logic: To create a counter with the points in each question and accumulate that in a total variable. Now I simply set the full threshold at the beginning to 75% and it gives the computed value automatically.
\newcommand{\pts}[1]{\if#11\textbf{[#1 point]}\else\textbf{[#1 points]}\fi\addtocounter{totalpts}{#1}}


\title{Homework 3}
%\author{Nicol\'as A Barnafi}
\date{}

\renewcommand{\vec}{\vectorsym}
\newcommand{\mat}{\matrixsym}
\newcommand{\ten}{\tensorsym}
\DeclareMathOperator{\grad}{\nabla}
\DeclareMathOperator{\dive}{\text{div}}
\DeclareMathOperator{\curl}{\text{curl}}
\DeclareMathOperator{\tr}{\text{tr}}
\DeclareMathOperator{\sym}{\text{sym}}
\newtheorem{remark}{Remark}
\newtheorem{definition}{Definition}
\newcommand{\R}{\mathbb{R}}
\newcommand{\D}{\mathcal{D}}

\newcommand{\tin}{\text{in}}
\newcommand{\ton}{\text{on}}

\newtheorem{theorem}{Theorem}
\newtheorem{lemma}{Lemma}

\begin{document}

\maketitle
 
\instruccion

\begin{enumerate}

    \item\pts{4}  Consider a bounded Lipschitz domain $\Omega\subset \R^d$, a function $\ten K:\Omega \to \R^{d\times d}$ in $L^\infty(\Omega,\R^{d\times d})$ such that it is symmetric, continuous, and positive definite almost everywhere ($\exists c, C>0: c|\vec x|^2 \leq \vec x^T\ten K\vec x \leq C|\vec x|^2$ a.e.), a function $\vec b$ in $L^\infty(\Omega, \R^d)$ such that $\dive {\vec b}=0$, $\vec b\cdot \vec n\geq 0$ on $\Gamma_N$ and $\vec b\in \vec H(\dive; \Omega)$, a positive scalar function $a>0$ in $L^\infty(\Omega)$, and an element $f$ of $[H^1(\Omega)]'$. In this context, consider the Advection-Diffusion-Reaction (ADR) problem
            $$ 
            \begin{aligned}
                -\dive \ten K\grad u + \vec b\cdot \grad u + a u &= f &&\text{ in $\Omega$}, \\
                u &= u_D &&\text{on $\Gamma_D$}, \\
                \ten K\grad u\cdot \vec n &= t &&\text{on $\Gamma_N$},
            \end{aligned}
            $$
            where $\partial \Omega=\overline\Gamma_D \cup \overline \Gamma_N$, $u_D$ is in $H^{1/2}(\Gamma_D)$, and $t$ in $H^{-1/2}(\Gamma_N)$. 
            \begin{enumerate}
                \item\pts{2} Find a weak formulation of this problem, clearly defining the solution space, the test function space, and the involved linear/bilinear forms. To do this, you only need to integrate by parts the second-order differential operator. Note that integrating the first-order term would generate new boundary terms that would be difficult to analyze.
                \item\pts{2} Prove that the found weak formulation has a unique solution using the Lax-Milgram Lemma. It will be useful to prove the following identity:
                   $$ \int_{\partial\Omega} u^2 (\vec b\cdot n)\,dS = \int_\Omega \dive (u\vec b) u\,dx + \int_\Omega (u\vec b)\cdot \grad u\,dx = \int_\Omega (\vec b\cdot \grad u + u \dive \vec b) u\,dx + \int_\Omega u\vec b\cdot \grad u\,dx $$
                   Clearly indicate where the hypotheses on the parameters are used.
            \end{enumerate}

        \item\pts{6} Consider a bounded Lipschitz domain $\Omega\subset \R^d$ and the following problem: Find $u_1, u_2:\Omega \to \R$ such that
            $$ 
            \begin{aligned}
                -\Delta u_1 + u_2 &= f_1  && \text{in $\Omega$} \\
                -\Delta u_2 - u_1 &= f_2  && \text{in $\Omega$},
            \end{aligned}
            $$
            given two functions $f_1,f_2$. 
            \begin{enumerate}
                \item\pts{2} Show, by integrating by parts each equation separately, what the appropriate boundary conditions for this problem are (Dirichlet, Neumann, or a mixture of them). 
                \item\pts{2} Consider homogeneous boundary conditions for both variables: 
                        $$ u_1 = u_2 = 0 \quad\text{on $\partial\Omega$}. $$
                       Write the weak formulation of the problem, clearly stating the spaces involved, and the required regularity for the functions $f_1, f_2$. For this, it will be useful to note that the solution space of your problem can be given by $V_0 = H_0^1(\Omega)\times H_0^1(\Omega)$, and therefore $\vec u \in V_0$ if and only if $\vec u = (u_1, u_2)$, for $u_1, u_2$ in $H_0^1(\Omega)$. 
                   \item\pts{2} Prove that the problem has a unique solution and write the a-priori stability bound that shows the continuity of the inverse. 
            \end{enumerate}

        \item\pts{10} Consider $\mu,\lambda>0$ (Lamé parameters), and a bounded Lipschitz domain $\Omega\subset\R^2$. Given a vector field $\vec u:\Omega\to \R^2$, which we will call displacement, we define its symmetric gradient as
            $$ \ten \varepsilon (\vec u) \coloneqq \frac 1 2\left(\grad \vec u + [\grad \vec u]^T\right), $$
            where the matrix $\grad \vec u$ is given pointwise by $[\grad \vec u]_{ij} = \frac{\partial u_i}{\partial x^j}$. The Hooke tensor is also defined as the following operator: 
            $$ \ten \sigma (\vec u) \coloneqq 2 \mu \ten \varepsilon(\vec u) + \lambda \dive (\vec u) \ten I. $$
            Given these definitions, a volumetric force $\vec f:\Omega\to \R^d$, and a boundary condition $\vec u_D$, the \emph{linear elasticity} problem is defined as
                $$ \begin{aligned}
                    -\dive \ten \sigma (\vec u) &= \vec f &&\text{ in $\Omega$},  \\
                    \vec u &= \vec u_D &&\text{on $\Gamma_D$},  \\
                    \ten \sigma \vec n &= \vec t &&\text{on $\Gamma_N$}. 
                \end{aligned} $$
                \begin{enumerate}
                    \item\pts{1} Prove that given a symmetric matrix $A$ and an antisymmetric matrix $B$, we have that 
                            $$ A : B = 0, $$
                            where $:$ is the matrix contraction or Frobenius product, given by $X:Y=\sum_{i,j} X_{ij}Y_{ij}$. 
                        \item\pts{1} Prove that the integration by parts formula in matrix form is given by
                        $$ \int_{\partial\Omega} \vec v\cdot (\ten \tau \vec n)\,dS = \int_\Omega \dive\ten \tau \cdot \vec v\,dx + \int_\Omega \ten\tau : \grad \vec v\,dx, $$
                        where $\ten \tau:\Omega \to \R^{d\times d}$ belongs to $\ten H(\dive; \Omega)\coloneqq[H(\dive;\Omega)]^d$, and $\vec v:\Omega \to \R^d$ belongs to $\vec H^1(\Omega)\coloneqq [H^1(\Omega)]^d$. Note that the $\dive$ operator applied to a matrix function $\ten \tau$ is applied row-wise, meaning it results in the following vector: 
                        $$ \dive \ten \tau = \begin{bmatrix} \dive \ten \tau_{1,\bullet} \\ \dive \ten \tau_{2,\bullet} \\ \vdots \\ \dive \ten \tau_{d, \bullet} \end{bmatrix} $$ 
                    \item\pts{2} Use the two previous results to show that the weak formulation of the elasticity problem is given by: Find $\vec u$ in $\vec V_{u_D}$ such that 
                        $$ \int_\Omega \ten \sigma(\vec u):\ten\varepsilon(\vec v)\,dx = \langle \vec f, \vec v\rangle_{V_0', V_0} + \langle \vec t, \vec v\rangle_{[H^{1/2}(\Gamma_N)]', H^{1/2}(\Gamma_N)} \qquad \forall \vec v \in \vec V_0,$$
                        where 
                        $$ \vec V_{u_D} = \{ \vec v \in \vec H^1(\Omega): \vec v = \vec u_D \text{ on $\Gamma_D$}\} $$
                        and
                        $$ \vec V_0 = \{ \vec v \in \vec H^1(\Omega): \vec v = \vec 0 \text{ on $\Gamma_D$}\}, $$
                        with the equalities in these definitions understood in the sense of traces. 
                    \item\pts{1} Show that the symmetric gradient $\varepsilon(\vec u)$ is such that 
                        $$ \| \ten \varepsilon(\vec u) \|_{0,\Omega} \leq C \| \grad \vec u \|_{0,\Omega} . $$
                    \item\pts{3} Prove Korn's inequality, given by
                        $$ \| \vec u \|_{1,\Omega} \leq C \|\ten \varepsilon(\vec u) \|_{0,\Omega} \qquad \forall \vec u \in [H_0^1(\Omega)]^d. $$
                            \emph{Hint: Extend Poincaré's proof for this case. Note that the condition $\ten \varepsilon(\vec u) = 0$ implies that $\vec u$ is a rigid body motion, i.e. a function of the form $\vec u = \left(\begin{smallmatrix} \alpha x_2 + \beta \\ -\alpha x_1 + \gamma \end{smallmatrix}\right)$, for arbitrary constants $\alpha, \beta, \gamma$. This holds only in 2D, so you can reduce the proof to the case $\Omega\subset \mathbb R^2$.}
                    \item\pts{2} Using the previous results, show that the linear elasticity problem has a unique solution if $|\Gamma_D|>0$ and show the stability bound.
                \end{enumerate}
\end{enumerate}

\todo[inline,color=white!90!black]{\textbf{Note: } We will open a forum on Canvas to review any typos and/or errors in the formulation.}
\end{document}
